\subsection{Fokusgruppe -- gjennomføringsplan og analyse}
\label{app:fokusgruppe}

\subsubsection*{Metodisk begrunnelse}

Fokusgrupper er egnet til å avdekke kollektive holdninger og meninger,
kartlegge kundemotivasjoner og teste aksept for nye produkter eller tjenester.
Metoden gir høy detaljgrad, men lav statistisk presisjon
\parencite{aspelund2026_forelesning7}. En svakhet ved fokusgrupper som metode
er at innsikten kan bli smal, at tolkning er subjektiv, og at enkeltpersoner
kan dominere diskusjonen og dermed påvirke andres svar.

\subsubsection*{Forskningsspørsmål}

Ziplines etableringsutfordring har to dimensjoner. Den første er kvalitativ og
handler om forbrukerholdninger og opplevd relevans. Den andre er strukturell
og handler om regulatoriske og infrastrukturelle barrierer. Fokusgruppen
adresserer de tre kvalitative forskningsspørsmålene. Det fjerde, strukturelle
spørsmålet er ikke egnet for fokusgrupper og belyses derfor gjennom
sekundærkilder.

\textbf{Kvalitative forskningsspørsmål.}
\begin{itemize}[leftmargin=1.5em, nosep]
  \item \textbf{FS1:} Oppfattes tjenesten som relevant sammenlignet med
  eksisterende alternativer?
  \item \textbf{FS2:} Hvilke holdninger har norske forbrukere til
  dronelevering, særlig med tanke på trygghet, personvern og støy?
  \item \textbf{FS3:} Vil bærekraft fungere som en verdiskapende faktor for
  dronelevering?
  \item \textbf{FS4:} Hva er norske forbrukeres pristerskel for dronelevering,
  og hvilke faktorer avgjør om en prispremie oppleves som akseptabel?
\end{itemize}

\textbf{Strukturelle forskningsspørsmål.}
\begin{itemize}[leftmargin=1.5em, nosep]
  \item \textbf{FS5:} Hvilke regulatoriske, infrastrukturelle og
  markedsmessige barrierer hindrer dronelevering i Norge?
\end{itemize}

\subsubsection*{Hypoteser}

Fire hypoteser ble formulert i forkant av gjennomføringen og evalueres mot
materialet i analysen lenger ned.

\begin{itemize}[leftmargin=1.5em, nosep]
  \item \textbf{H1} (knyttet til FS1): Forbrukere opplever ikke leveringstid
  som frustrerende nok til at 15-minutters dronelevering oppfattes som en
  løsning på et reelt og uløst behov.
  \item \textbf{H2} (knyttet til FS2): Problemer knyttet til trygghet,
  personvern og støy veier tyngre enn interessen for hastighet.
  \item \textbf{H3} (knyttet til FS3): Bærekraft nevnes som positivt, men er
  ikke en reell driver for valg av leveringstjeneste.
  \item \textbf{H4} (knyttet til FS4): Forbrukerne har en lav pristerskel for
  dronelevering og vil kun akseptere en prispremie i situasjoner der
  tjenesten løser et behov som eksisterende alternativer ikke dekker.
\end{itemize}

\subsubsection*{Utvalg og rekruttering}

Fokusgruppen bestod av fire deltakere fra samme familie og utgjør dermed en
homogen gruppe som ikke er representativ for den generelle befolkningen
eller ulike kundesegmenter. Deltakerne er i tillegg lite aktive brukere av
hjemlevering, noe som begrenser overføringsverdien til hyppige brukere.
Anbefalt størrelse for en fokusgruppe er 4--12 deltakere med én eller flere
moderatorer \parencite{aspelund2026_forelesning7}, og gruppen ligger dermed
i nedre sjikt av det anbefalte intervallet. Færre deltakere reduserer
bredden i meninger som fanges opp, men gir samtidig den enkelte mer tid til
å utdype sine synspunkter. Det ble kun gjennomført én fokusgruppe, noe som
gir et smalt grunnlag der enkeltstemmer kan få uforholdsmessig stor vekt.
Funnene må derfor tolkes med forbehold og ses i sammenheng med de
kvantitative resultatene fra spørreundersøkelsen
(vedlegg~\ref{app:sporreundersokelse}).

\subsubsection*{Praktisk oppsett}

Fokusgruppen ble gjennomført rundt et kjøkkenbord med en varighet på omtrent
60 minutter, og plasseringen sikret at samtlige deltakere kunne se
hverandre. Moderator ledet samtalen og stilte spørsmålene uten å styre
svarene, mens observatøren noterte hvem som sa hva, hvilke temaer som
utløste engasjement eller motstand, samt relevant kroppsspråk og reaksjoner.
Moderator- og observatørrollen ble utført av samme person, noe som i praksis
gjorde det krevende å styre samtalen og notere observasjoner samtidig.

\subsubsection*{Gjennomføring av fokusgruppe}

Samtalen var strukturert fra bredt til konkret, slik at Zipline-konseptet
først ble introdusert etter at deltakernes generelle holdninger og vaner
knyttet til hjemlevering var kartlagt. Dette ble gjort for å redusere
risikoen for at konseptet styrte de innledende svarene. Fokusgruppen ble
gjennomført anonymt, og deltakerne omtales som D1, D2, D3 og D4 gjennom
resten av analysen.

\paragraph{Del A -- Oppvarming (5 min).}\mbox{}\\

\textit{Hva spiste du sist gang du bestilte mat hjem?}
\begin{itemize}[leftmargin=1.5em, nosep]
  \item \textbf{D1:} Sushi.
  \item \textbf{D2:} Sushi.
  \item \textbf{D3:} Sushi.
  \item \textbf{D4:} Sushi.
\end{itemize}

\textit{Hvor ofte bestiller dere hjemlevering av mat, og hva er den typiske
situasjonen når det skjer?}
\begin{itemize}[leftmargin=1.5em, nosep]
  \item \textbf{D1:} Sjeldent, men skjer gjerne i forbindelse med møter
  knyttet til verv på skolen.
  \item \textbf{D2:} Mindre enn én gang i året. Bestiller i så fall når det
  skjer noe sosialt med venner.
  \item \textbf{D3:} Bestilt hjemlevering veldig sjeldent.
  \item \textbf{D4:} Bestiller hovedsakelig store mengder pizza til
  arrangementer man er med på å arrangere.
\end{itemize}

\paragraph{Del B -- Problem og behov (10 min).}\mbox{}\\

\textit{Hva er det egentlige problemet du løser når du bestiller hjemlevering
av mat? Handler det om tid, matlyst, energi, noe annet?}
\begin{itemize}[leftmargin=1.5em, nosep]
  \item \textbf{D1:} Det handler primært om å spare tid, og at det er enkelt
  og praktisk når man er lat eller fyllesyk og ikke orker å lage mat selv.
  \item \textbf{D2:} Det handler om å spare tid. Man har heller ikke alltid
  mulighet til å hente mat selv uten tilgang på bil.
  \item \textbf{D3:} Ikke besvart.
  \item \textbf{D4:} Det handler primært om å spare tid og forenkle
  logistikken rundt mat ved arrangementer.
\end{itemize}

\textit{Hva fungerer bra med tjenestene du har brukt? Hva er irriterende
eller mangelfullt?}
\begin{itemize}[leftmargin=1.5em, nosep]
  \item \textbf{D1:} Tjenestene er ikke alltid helt pålitelige når det
  gjelder ventetid.
  \item \textbf{D2:} Det er enkelt å bestille, og praktisk å logge inn i en
  app, skrive hva man vil ha og hvor det skal leveres.
  \item \textbf{D3:} Det er bredt utvalg å velge mellom.
  \item \textbf{D4:} Det er positivt at man får med bestikk og lignende, og
  at de tilbyr ekstra produkter som rømmedressing. I tillegg er
  betalingsprosessen enkel og grei.
\end{itemize}

\textit{Hvor viktig er pris, hastighet, utvalg, pålitelighet og enkelhet for
deg? Ranger alle fra 1 (lavt) til 5 (høyt).}

\begin{table}[H]
\centering
\caption{Deltakernes rangering av leveringsattributter (1--5)}
\label{tab:fokusgruppe_rangering}
\renewcommand{\arraystretch}{1.2}
\begin{tabular}{@{} l c c c c c @{}}
\toprule
 & \textbf{Pris} & \textbf{Hastighet} & \textbf{Utvalg} & \textbf{Pålitelighet} & \textbf{Enkelhet} \\
\midrule
D1 & 4 & 3 & 2 & 3 & 3 \\
D2 & 5 & 3 & 2 & 5 & 3 \\
D3 & 2 & 5 & 5 & 5 & 5 \\
D4 & 1 & 3 & 4 & 3 & 4 \\
\bottomrule
\end{tabular}
\end{table}

\paragraph{Del C -- Kundereisen (10 min).}\mbox{}\\

\textit{Når du bestemmer deg for å bestille mat, hva skjer fra du kjenner
behovet til du faktisk trykker bestill?}
\begin{itemize}[leftmargin=1.5em, nosep]
  \item \textbf{D1:} Først vurderer D1 hva et tilsvarende måltid ville kostet
  på Rema 1000, og regner ut differansen mot hva det ville kostet hos
  Foodora/Wolt. Deretter spør D1 seg selv om tiden man sparer på å slippe å
  lage maten er verdt pengene.
  \item \textbf{D2:} Starter med å tenke på hva slags mat man har lyst på,
  søker deretter opp restauranter som tilbyr det, og blar gjennom
  alternativene før man tar en beslutning.
  \item \textbf{D3:} Spør seg selv hvor det er enkelt å bestille mat, hvem
  som har god oversikt over meny og enkel betaling. Det er også en fordel om
  det er et sted man har bestilt fra tidligere og er kjent med.
  \item \textbf{D4:} Beslutningen kan tas alt fra dagen før til rett før man
  trenger maten.
\end{itemize}

\textit{Hvem eller hva påvirker valget ditt, eksempelvis anmeldelser, venner,
pris, tidligere erfaring?}
\begin{itemize}[leftmargin=1.5em, nosep]
  \item \textbf{D1:} Valget påvirkes hovedsakelig av anbefalinger fra venner,
  gode tilbud og tidligere erfaring med stedet.
  \item \textbf{D2:} De viktigste faktorene er anbefaling fra venner,
  kjennskap til stedet, gode tilbud, kampanjer og studentrabatt.
  \item \textbf{D3:} Valget påvirkes av anbefaling fra venner, reklame og at
  restauranten man bestiller fra ser innbydende ut.
  \item \textbf{D4:} Valget avhenger sterkt av tidligere erfaring med
  restauranter.
\end{itemize}

\textit{Hva er det som til slutt avgjør om du bestiller eller ikke?}
\begin{itemize}[leftmargin=1.5em, nosep]
  \item \textbf{D1:} At smaken og tidsbesparing gjør prisen verdt det.
  \item \textbf{D2:} At man har lyst på maten og at prisen oppleves som
  akseptabel.
  \item \textbf{D3:} At man finner det man ønsker og er overbevist om at
  maten er god.
  \item \textbf{D4:} Om alle som skal ha mat faktisk møter opp.
\end{itemize}

\paragraph{Del D -- Test av Zipline-konsept (25 min).}\mbox{}\\
Deltakerne fikk følgende beskrivelse opplest på en nøytral måte: «Zipline er
et selskap som leverer mat via drone. Dronen flyr 100 meter over bakken og
senker ned pakken uten å lande. Leveringstider ca. lik som for annen
hjemlevering. Tenk dere at dette fantes i din by.»

\textbf{Førsteinntrykk (FS1).}

\textit{Hva er den aller første tanken din?}
\begin{itemize}[leftmargin=1.5em, nosep]
  \item \textbf{D1:} Konseptet virker lovende dersom det er billigere,
  raskere og sikkert. Samtidig er det bekymring knyttet til støy og følelsen
  av å bli overvåket.
  \item \textbf{D2:} Den første reaksjonen er skepsis knyttet til
  leveringssikkerhet og hvordan man sikrer at maten kommer frem til riktig
  person og ikke blir stjålet. Det er også usikkerhet rundt hvordan
  tjenesten håndterer dårlig vær.
  \item \textbf{D3:} Konseptet virker nyttig for folk som bor
  avsidesliggende til, for eksempel på den andre siden av en fjord, eller er
  på hytta. Det påpekes også at mange forbinder droner med krig, og at folk
  derfor kan synes de er litt ubehagelige. Det er også ubehag knyttet til at
  dronene flyr rundt og filmer folk med kamera, selv om det ikke er det som
  er meningen.
  \item \textbf{D4:} Dette oppleves som naturlig teknologisk utvikling og at
  det er den veien verden beveger seg i.
\end{itemize}

\textit{Hva virker attraktivt?}
\begin{itemize}[leftmargin=1.5em, nosep]
  \item \textbf{D1:} Det er ikke viktig hvordan maten blir levert så lenge
  det er billigere og kjappere enn andre alternativer.
  \item \textbf{D2:} Konseptet virker nyskapende og spennende. Deltakeren er
  nysgjerrig og motivert til å prøve tjenesten fordi det er noe nytt og
  spennende.
  \item \textbf{D3:} Rask levering og bedre hygiene trekkes frem som
  attraktivt, ettersom tradisjonelle leveringsmetoder som sykkelbud og
  kjøretøy kan oppleves som uhygieniske.
  \item \textbf{D4:} Miljøprofilen virker attraktiv. Det ville også vært
  attraktivt hvis dronen kunne hentet opp mat fra flere ulike steder i én og
  samme leveranse.
\end{itemize}

\textit{Hva virker uklart eller usikkert?}
\begin{itemize}[leftmargin=1.5em, nosep]
  \item \textbf{D1:} Det er usikkerhet rundt overleveringen av maten når man
  ikke møter et menneske. Det er uklart hvordan mottaksprosessen fungerer.
  \item \textbf{D2:} Værforhold trekkes frem som en bekymring. Det er
  usikkerhet rundt pålitelighet, om dronen klarer å levere til riktig person.
  Det er også usikkerhet om dronen kan håndtere store bestillinger, holde
  maten varm og transportere den stabilt uten at maten velter.
  \item \textbf{D3:} Det er bekymring rundt værforhold, og om tjenesten er
  pålitelig nok i dårlig vær.
  \item \textbf{D4:} Det er usikkerhet rundt om dronen klarer å gjennomføre
  leveranser under krevende værforhold som vind, regn og snø.
\end{itemize}

\textbf{Barrierer og drivere (FS2).}

\textit{Hva skal til for at du tester tjenesten, og hva ville stoppet deg?}
\begin{itemize}[leftmargin=1.5em, nosep]
  \item \textbf{D1:} Tjenesten måtte være billig og bevist at den fungerer,
  og helst bli en norm folk flest bruker. Det som ville stoppet valget er at
  det er dyrt og ikke trygt.
  \item \textbf{D2:} Det viktigste er at man kan stole på at tjenesten er
  pålitelig og at de har et godt tilbud. Et samarbeid med influencere kunne
  gitt et nyttig innblikk i hvordan tjenesten fungerer. Det ville også
  hjulpet dersom de samarbeidet med store restaurantkjeder som Peppes og
  Sumo, som man er fortrolig med, stoler på og har god erfaring med.
  \item \textbf{D3:} Anbefaling fra noen man kjenner, nysgjerrighet på selve
  leveringsmåten, eller bekjentskap med tjenesten gjennom reklame er viktig.
  Det ville også hjulpet dersom de samarbeidet med store aktører man er
  fortrolig med og stoler på.
  \item \textbf{D4:} Det som ville stoppet valget er usikkerhet rundt
  leveringen, forsinkelser, om maten holder seg varm og om den kommer
  tidsnok.
\end{itemize}

\textit{Ville du vært komfortabel med en drone over nabolaget ditt?}
\begin{itemize}[leftmargin=1.5em, nosep]
  \item \textbf{D1:} Det er i utgangspunktet greit, men det er viktig at man
  vet om det er en matdrone eller en drone som filmer. Usikkerhet rundt
  dette fører til at man føler at man har mindre privatliv.
  \item \textbf{D2:} Det ville ikke vært komfortabelt hvis det ble for mange
  droner. Én til to droner er greit, men flere droner hver eneste dag er
  ikke akseptabelt. Det ville vært lettere å akseptere utenfor sentrum hvor
  det er mindre bebyggelse.
  \item \textbf{D3:} Droner forbindes med krig, og det ville ikke vært så
  komfortabelt. Det ville også vært ubehagelig hvis naboene bestilte mat og
  man ikke visste at det kom en drone. Det ville vært lettere å akseptere
  utenfor sentrum hvor det er mindre bebyggelse.
  \item \textbf{D4:} Det er ikke særlig ønskelig med droner over hus og
  nabolag hele tiden.
\end{itemize}

\textit{Hva tenker du om støy og personvern?}
\begin{itemize}[leftmargin=1.5em, nosep]
  \item \textbf{D1:} Det er irriterende hvis dronen bråker mye. Det går i
  utgangspunktet fint, men hvis det er mistanke om at dronen brukes til noe
  annet enn matlevering er det ikke akseptabelt. Det er derfor viktig at
  dronene er godt merket slik at man kan se hvem som eier dem.
  \item \textbf{D2:} Droner oppleves som forstyrrende for personvernet. Hvor
  irriterende støyen er avhenger av mengden. Dronene bør være godt merket
  slik at man vet at det dreier seg om matlevering og ikke noe annet.
  \item \textbf{D3:} Det er enighet med D2 om at droner bør være godt merket,
  ettersom tydelig merking gjør at man føler seg tryggere.
  \item \textbf{D4:} Det er bekymring knyttet til mye støy over huset, og
  droner oppleves heller ikke som særlig gunstig for personvernet.
\end{itemize}

\textbf{Bærekraft (FS3).}

\textit{Zipline hevder droneleveranser gir 97\,\% lavere
CO\textsubscript{2}-utslipp enn tradisjonell levering. Ville det påvirket
valget ditt i praksis, eller er det mer en bonus?}
\begin{itemize}[leftmargin=1.5em, nosep]
  \item \textbf{D1:} Miljøprofilen ville påvirket valget positivt, men prisen
  ville uansett trumfet hensynet til bærekraft. Det kan tippe valget i en
  viss retning, men er ikke en av de primære faktorene.
  \item \textbf{D2:} Miljøprofilen hadde påvirket valget dersom kostnadene
  var like som for eksisterende alternativer.
  \item \textbf{D3:} Miljøprofilen hadde påvirket valget dersom kostnadene
  var like som for eksisterende alternativer.
  \item \textbf{D4:} Miljøprofilen oppleves mer som en bonus og ville ikke
  påvirket valget i praksis.
\end{itemize}

\textbf{Situasjonsbruk.}

\textit{Kan du beskrive en konkret situasjon der Zipline hadde passet
perfekt?}
\begin{itemize}[leftmargin=1.5em, nosep]
  \item \textbf{D1:} En fjelltur der man kan få levert mat direkte til en
  GPS-posisjon.
  \item \textbf{D2:} Tjenesten hadde passet godt for folk som bor usentralt
  og kronglete til.
  \item \textbf{D3:} Hyttetur og båttur trekkes frem som situasjoner der
  dronelevering hadde vært nyttig.
  \item \textbf{D4:} En fjelltur der man kan få levert mat direkte til en
  GPS-posisjon.
\end{itemize}

\textit{Er det situasjoner der du aldri ville valgt det?}
\begin{itemize}[leftmargin=1.5em, nosep]
  \item \textbf{D1:} Hvis man har en nabo som ikke liker droner, ville det
  ikke vært aktuelt å bruke tjenesten.
  \item \textbf{D2:} Ville unngått å bruke tjenesten i byen, med unntak av
  situasjoner der man for eksempel har brukket foten og ikke kommer seg ut.
  Det ville heller ikke vært aktuelt dersom man bor i nærheten av en
  flyplass eller militærleir.
  \item \textbf{D3:} I tettbebygde strøk og i blokk ville tjenesten ikke
  blitt valgt.
  \item \textbf{D4:} Dersom man har en nabo som er svært negativ til droner
  ville det ikke vært aktuelt å bruke tjenesten.
\end{itemize}

\textbf{Segment og posisjonering.}

\textit{Hvem tror du dette passer best for? Beskriv den typiske
Zipline-brukeren.}
\begin{itemize}[leftmargin=1.5em, nosep]
  \item \textbf{D1:} Den typiske brukeren beskrives som unge folk som er
  tidlig ute med å ta i bruk ny teknologi.
  \item \textbf{D2:} Den typiske brukeren beskrives som en ung mann rundt 30
  år uten kjæreste, som bor usentralt og er teknologientusiast. Det pekes
  også på at tjenesten passer for folk som ikke har tilgang på Foodora eller
  lignende.
  \item \textbf{D3:} Den typiske brukeren beskrives som en gjeng med menn
  mellom 20 og 60 år som synes droner og teknologi er kult og spennende.
  \item \textbf{D4:} Den typiske brukeren beskrives som unge mennesker i
  20-årene og begynnelsen av 30-årene som har begynt i jobb og har penger
  til å bestille mat, og som er vant til å bruke teknologi i hverdagen.
\end{itemize}

\paragraph{Del E -- Prioritering og avslutning (10 min).}\mbox{}\\

\textit{Hva er de viktigste tingene Zipline må få på plass for at du skal
velge dem?}
\begin{itemize}[leftmargin=1.5em, nosep]
  \item \textbf{D1:} Tjenesten må være pålitelig, billig og trygg.
  \item \textbf{D2:} Pålitelighet er avgjørende, i tillegg til at
  bestillingsmåten må være profesjonell, oversiktlig og enkel.
  \item \textbf{D3:} Det må være på plass et system som gjør leveringen og
  det å ta imot maten så enkelt som mulig.
  \item \textbf{D4:} Pålitelighet trekkes frem som det viktigste.
\end{itemize}

\textit{Hva må være på plass for at dette skal føles trygt?}
\begin{itemize}[leftmargin=1.5em, nosep]
  \item \textbf{D1:} Det er viktig at man ikke opplever at teknologien
  misbrukes til andre formål, og at det ikke skjer uhell.
  \item \textbf{D2:} Det må være kjent at selskapet har fått de nødvendige
  tillatelsene. Det oppleves også som tryggere dersom det er mennesker som
  styrer dronene, og ikke roboter.
  \item \textbf{D3:} Det er avgjørende at leveringsmåten er godkjent og
  lovlig. I tillegg bør tjenesten reklameres godt på forhånd slik at det
  blir allment kjent hva som skjer når noen får levert mat på denne måten.
  \item \textbf{D4:} Det må være bevist at dronene fungerer bra og ikke
  krasjer. Teknologien må være gjennomprøvd før man kan føle seg trygg på å
  bruke tjenesten.
\end{itemize}

\subsubsection*{Analyse av fokusgruppe -- Zipline i det norske markedet}

\textbf{FS1 -- Opplevd relevans sammenlignet med eksisterende alternativer.}
Deltakernes primære motivasjon for hjemlevering var tidsbesparelse og
praktisk logistikk, ikke hastighet i seg selv, og D1 og D4 bestilte
hovedsakelig i situasjoner der 15-minutters levering var irrelevant.
Tjenesten ble først oppfattet som relevant i rurale kontekster som
fjellturer, hytter og båtturer, mens urbane hverdagssituasjoner manglet et
tydelig problem å løse. H1 støttes dermed av materialet.

\textbf{FS2 -- Holdninger knyttet til trygghet, personvern og støy.}
Tre av fire deltakere trakk spontant frem værutsatthet, leveringssikkerhet og
personvern som sentrale bekymringer, og D3 assosierte droner med krig og
overvåking. Samtlige understreket at tydelig merking var en forutsetning for
legitimitet, og pålitelighet og trygghet dominerte over hastighet da
deltakerne rangerte hva som skulle til for å velge tjenesten. H2 støttes
dermed av materialet.

\textbf{FS3 -- Bærekraft som verdiskapende faktor.}
Bærekraft ble ikke spontant nevnt som motivasjon, og først når 97\,\% lavere
CO\textsubscript{2}-utslipp ble presentert kom det respons. D2 og D3 ville
latt miljøprofilen påvirke valget kun dersom prisen var lik eksisterende
alternativer, D1 omtalte det som noe som kunne tippe valget uten å være
primært, og D4 opplevde det kun som en bonus. H3 støttes dermed av
materialet.

\textbf{FS4 -- Pristerskel og akseptabel prispremie.}
Prissensitiviteten var gjennomgående høy, og D1 og D2 rangerte pris høyest
eller nest høyest av alle attributter. Ingen av deltakerne nevnte
hverdagssituasjoner der de ville akseptert en prispremie for hastighet
alene, men flere antydet at en konkret bruksverdi som GPS-levering på fjellet
kunne endre regnestykket. H4 støttes dermed av materialet.

\textbf{Samlet vurdering.}
Fokusgruppen avdekker en konsistent prioriteringsrekke der pålitelighet, pris
og trygghet veier langt tyngre enn hastighet og bærekraft. Et gjennomgående
mønster er at Zipline oppleves som mest relevant i rurale bruksarenaer, mens
urbane hverdagssituasjoner mangler et sterkt nok følt behov til å
rettferdiggjøre tjenesten. Deltakerne peker samtidig på tydelig merking,
samarbeid med etablerte aktører og gjennomprøvd teknologi som forutsetninger
for at terskelen for førstegangsbruk skal senkes.
